\documentclass[10pt,a4paper]{amsart}
\usepackage[margin=19mm]{geometry}
\usepackage{amsmath,amssymb,mathtools}
\usepackage[scr=rsfs,cal=euler]{mathalfa}
\usepackage{xfrac}
\usepackage[unicode,hidelinks,bookmarks=false]{hyperref}
\newcommand{\ie}{i.e.\ }
\newcommand{\cf}{cf.\ }
\newcommand{\resp}{respectively\ }
\newcommand{\Subsection}{\subsection}
\providecommand{\nobreakdash}{\nobreak\hbox{-}}
\setlength{\emergencystretch}{2em}
\multlinegap0pt
\usepackage{float}
\usepackage{amssymb}
\usepackage{enumerate}
\usepackage{xcolor}
\newtheorem{corollary}{Corollary}
\newtheorem{theorem}{Theorem}
\title{Secure coalitions in graphs}
\author{Swathi Shetty, Sayinath Udupa N. V., B. R. Rakshith}
\date{Source: arXiv:2511.21170v2 (CC BY 4.0)}
\begin{document}
\maketitle

\noindent\emph{Source and licence.} Secure coalitions in graphs. Version arXiv:2511.21170v2 (CC BY 4.0), distributed by arXiv under the Creative Commons Attribution 4.0 International licence, http://creativecommons.org/licenses/by/4.0/, as recorded in the arXiv licence metadata for this exact version and shown on the abstract page https://arxiv.org/abs/2511.21170v2. Full licence notice: \texttt{licenses/arxiv-2511.21170-CC-BY-4.0.txt}.\par\smallskip

\noindent\textit{Editorial scope.} Counted unit: the article's theorem treen-2 (source0.tex lines 398--401). The article's complete original proof of it is delivered with it (source0.tex lines 402--526, one \texttt{proof} environment of 125 lines). No mathematics is written, repaired or re-derived: the statement and the proof are the article's own lines, in the article's own notation, and no retained line is substituted or re-flowed.  Retained uncounted as the source-local context the proof needs: lines 392--394.  One declared adaptation: exactly 1 ancillary strings inside the retained spans are omitted (image-inclusion commands and, where the source repeats a label, the repeated label definitions) and no other line differs from the source; the figure environments, with their placement, captions and labels, are kept, and the package records the omitted strings in fidelity receipts bound to the affected bodies.  No retained line contains a citation, so the wrapper carries no bibliography and no bibliography entry.  The retained lines contain the article's own diagram code, which the wrapper typesets with the standard tikz packages; no image file is delivered and the package declares no asset dependency.  Editorial formatting only, in addition: an AMS \texttt{amsart} preamble that restates the article's own declarations and macros used by the retained lines, namely the environments \texttt{corollary}, \texttt{theorem} on separate counters, together with the standard packages those lines need.  The retained environments print in this document as Corollary 1, Theorem 1 in order of appearance, whereas the article prints the same items with the numbering of its own full document.  The complete native source of the article as distributed in the arXiv e-print for this exact version is bundled unchanged as \texttt{sources/189834/source0.tex}.
\begin{corollary}\label{nsecnumber1}
For any tree $T$ of order $n$, $SEC(T)=n$ if and only if $T\cong P_n$, where $n\le 4$.
\end{corollary}

\begin{theorem}\label{treen-2}
Let $T$ be a tree of order $n\ge 5$. If $T\not\cong P_5$, then
$SEC(T)\le n-2$.
\end{theorem}

\begin{proof}
By Corollary~\ref{nsecnumber1}, we have $SEC(T)<n$ for every tree $T$ of order $n\ge 5$. Therefore, it suffices to show that, for $n\ge 5$ and $T\not\cong P_5$, we have $SEC(T)\neq n-1$.

Suppose, to the contrary, that $SEC(T)=n-1$. Let $\pi$ be a secure coalition partition of $T$ with $|\pi|=n-1$. Then exactly one set of $\pi$ has cardinality $2$, say $A$, while remaining are singleton sets.

Let $x\in V(T)$ be a pendant vertex such that $x\notin A$.

\textbf{Claim 1.} The singleton set $\{x\}$ cannot form a secure coalition with any singleton set of $\pi$.

Suppose, to the contrary, that $\{x\}$ forms a secure coalition with a singleton set $\{z\}\in\pi$. Then $S=\{x\}\cup\{z\}$
is a secure dominating set of $T$. Since $x$ is a pendant vertex, the vertex $z$ must dominate every vertex not dominated by $x$. Moreover, because $T$ is a tree, no two vertices of $N(z)$ are adjacent; otherwise, they would form a cycle together with $z$.
Now let $w\in N(z)\setminus N[x]$. Since $w$ is dominated by $z$, the set $(S\setminus\{z\})\cup\{w\} =\{x,w\}$
must be a dominating set of $T$. However, one can easily observe that, the set $\{x,w\}$ cannot dominate all vertices of $T$. This contradicts the assumption that $S$ is a secure dominating set. Therefore, $\{x\}$ cannot form a secure coalition with any singleton set of $\pi$.

Hence, the only possible secure coalition partner of $\{x\}$ is the set $A$. Since every tree has at least two pendant vertices, let $y\neq x$ be another pendant vertex of $T$. Then $A$ must contain either $y$ or the support vertex adjacent to $y$; otherwise, $A\cup\{x\}$ fails to dominate $y$ and therefore is not a dominating set.

We now consider the following cases.

\textbf{Case 1.} The vertex $y$ belongs to $A$.

Let $A=\{a,y\}$, where $a\neq y$. Since $A\cup\{x\}$ is a dominating set, the vertex $a$ must dominate all vertices that are not dominated by $x$ and $y$.

\textbf{Subcase 1.1.} The vertex $a$ is adjacent to neither $x$ nor $y$.

\textbf{Claim 1.1.1.} The vertex $a$ is not a pendant vertex.

Suppose, to the contrary, that $a$ is a pendant vertex. Since $\{x\}\cup A=\{x,a,y\}$ is a dominating set, it follows that $T$ must be isomorphic to one of the graphs $G_1$ or $G_2$ shown in Fig.~\ref{figseccoalition}.

\begin{figure}[H]
\centering

\caption{Graphs $G_1$ and $G_2$.}
\label{figseccoalition}
\end{figure}

In both $G_1$ and $G_2$, the singleton set corresponding to the support vertex $y_1$ of $y$ cannot form a secure coalition with any part of $\pi$. This contradicts the assumption that $\pi$ is a secure coalition partition. Hence, $a$ is not a pendant vertex.

Furthermore, since $T\not\cong P_5$, every vertex adjacent to $a$ that does not belong to $N(x)\cup N(y)$ must be a pendant vertex. Let $a_1$ be such a pendant vertex adjacent to $a$.

\textbf{Claim 1.1.2.} The singleton set $\left\{a_1\right\}$ does not form a secure coalition with any set of $\pi$.

Since $a_1$ is a pendant vertex, an argument identical to that used in Claim~1 shows that $\left\{a_1\right\}$ cannot form a secure coalition with any singleton set of $\pi$. Therefore, its only possible secure coalition partner is the set $A=\left\{a,y\right\}$.

However, $S=\{a_1\}\cup A=\left\{a_1,a,y\right\}$
is not a dominating set, since none of its vertices dominates the pendant vertex $x$. This contradiction shows that $\left\{a_1\right\}$ cannot form a secure coalition with any set of $\pi$. Hence, the claim follows.

\textbf{Subcase 1.2.} The vertex $a$ is adjacent to exactly one of $x$ and $y$.

\textbf{Claim 1.2.1.} The singleton set corresponding to the support vertex $y_1$ of $y$ cannot form a secure coalition with any set of $\pi$.

First suppose that $a\sim y$ and $a\not\sim x$. Then $a=y_1$. Since $n\ge 5$, the vertex $y_1$ is adjacent to a pendant vertex other than $y$, say $a_2$. By Claim~1, the singleton set $\left\{a_2\right\}$ cannot form a secure coalition with any singleton set of $\pi$. Moreover, its only possible coalition partner is $A=\left\{y,y_1\right\}$, but $\left\{a_2,y,y_1\right\}$ is not a secure dominating set. A contradiction.

Now suppose that $a\sim x$ and $a\not\sim y$. Consider the singleton set $\left\{y_1\right\}$. Since neither $y_1$ nor any vertex other than $a$ dominates both pendant vertices $x$ and $a_2$, the set $\left\{y_1\right\}$ cannot form a secure coalition with any singleton set of $\pi$. Furthermore,
$\left\{y_1\right\}\cup A=\left\{y_1,a,y\right\}$ is not a secure dominating set. Hence, $\left\{y_1\right\}$ does not form a secure coalition with any set of $\pi$, a contradiction.

Therefore, in either case, the singleton set corresponding to the support vertex $y_1$ of $y$ cannot form a secure coalition with any part of $\pi$. This proves the claim.

\textbf{Subcase 1.3.} The vertex $a$ is adjacent to both $x$ and $y$.

Since $x$ and $y$ are pendant vertices and every neighbor of $a$ other than those in $N(x)\cup N(y)$ is also a pendant vertex, it follows that $T\cong S_n$.
Observe that $\left\{x\right\}\cup A=\left\{x,a,y\right\}$
is a dominating set of $T$. Since $n\ge 5$, there exists a pendant vertex $a_1\neq x,y$ adjacent to $a$. As $a_1$ is dominated only by $a$, replacing $a$ with $a_1$ yields
$(\left\{x,a,y\right\}\setminus\left\{a\right\})\cup\left\{a_1\right\}=\left\{x,y,a_1\right\}$, which is not a dominating set of $T$. Therefore, $\left\{x,a,y\right\}$ is not a secure dominating set. Consequently, $\left\{x\right\}$ and $A$ do not form a secure coalition, a contradiction.

\textbf{Case 2.} The support vertex $y_1$ of $y$ belongs to $A$, and there exists a vertex $a\neq y$ such that $A=\{y_1,a\}$.

\textbf{Subcase 2.1.} The vertex $x$ is adjacent to $y_1$.

\textbf{Claim 2.1.1.} Every vertex adjacent to $y_1$ is a pendant vertex.

Suppose, to the contrary, that there exists a vertex $y_1'\in N(y_1)$ with $d(y_1')\ge 2$. Since $T$ is a tree, $y_1'$ is adjacent to a pendant vertex, say $y_2'$.

First assume that $a=y_1'$. Then the singleton set $\{y_2'\}$ cannot form a secure coalition with any singleton set of $\pi$. Moreover,
$\left\{y_2'\right\}\cup\left\{y_1,y_1'\right\}$
is a dominating set but not a secure dominating set, yielding a contradiction.

Next assume that $a=y_2'$. By a similar argument,
$\left\{y_1'\right\}\cup\left\{y_1,y_2'\right\}$
is a dominating set but not a secure dominating set, again a contradiction.

Finally, suppose that $a\notin\left\{y_1',y_2'\right\}$. Since $a$ cannot be adjacent to both $x$ and $y$, it follows that
$\left\{y_1'\right\}\cup\left\{y_1,a\right\}$ is not a secure dominating set, a contradiction. Hence, every neighbor of $y_1$ is a pendant vertex. Consequently, $T$ is a star, that is,
$T\cong S_n$. However, for $n\ge 5$,
$SEC(S_n)=3<n-1$, contradicting the assumption that $SEC(T)=n-1$. 

\textbf{Subcase 2.2.} The vertex $x$ is not adjacent to $y_1$.

\textbf{Claim 2.2.1.} The vertex $a$ is adjacent to $x$.

Suppose, to the contrary, that $a\not\sim x$. Then the singleton set $\{y\}$ cannot form a secure coalition with any singleton set of $\pi$. Moreover, $\{y\}\cup\{a,y_1\}$
fails to be a dominating set. This contradiction proves that $a\sim x$.

\textbf{Claim 2.2.2.} The vertex $y_1$ is not adjacent to any pendant vertex other than $y$.

Suppose that $y_1$ is adjacent to a pendant vertex $y_3\neq y$. Then $\{x\}\cup\{y_1,a\}$
is not a secure dominating set, contradicting the fact that $\{x\}$ and $A=\{y_1,a\}$ must form a secure coalition. Hence, $y_1$ is not adjacent to any pendant vertex other than $y$.

\textbf{Claim 2.2.3.} The vertex $a$ is not adjacent to any pendant vertex other than $x$.

Suppose that $a$ is adjacent to a pendant vertex $x_2\neq x$. Then the singleton set $\{y\}$ cannot form a secure coalition with any set of $\pi$, a contradiction. Therefore, $a$ is not adjacent to any pendant vertex other than $x$.

\textbf{Claim 2.2.4.} The vertices $y_1$ and $a$ are adjacent.

Suppose that $y_1\not\sim a$. Then the singleton set $\{y\}$ cannot form a secure coalition with any set of $\pi$, contradicting the assumption that $\pi$ is a secure coalition partition. Hence, $y_1\sim a$.

Furthermore, one can verify that if either $a$ or $y_1$ is adjacent to a non-pendant vertex, then the resulting tree cannot satisfy $SEC(T)=n-1$. Therefore, this subcase also leads to a contradiction.\\
Furthermore, Claims~2.2.1--2.2.4 imply that every neighbor of $a$ other than $x$ and every neighbor of $y_1$ other than $y$ is a pendant vertex. Consequently, the resulting structure does not admit a secure coalition partition of cardinality $n-1$, contradicting the assumption that $SEC(T)=n-1$.\\[2mm]
We now consider the cases in which $x\in A$.

\textbf{Case 3.} The vertex $y\notin A$.

This case is analogous to Case~1. Indeed, by interchanging the roles of $x$ and $y$ in the arguments of Case~1, we obtain the  contradiction. Hence, this case cannot occur.

\textbf{Case 4.} The vertex $y\in A$.

Since $n\ge 5$ and $T\not\cong P_5$, there exists a vertex $z\in V(T)$ that is not dominated by either $x$ or $y$. Consequently, every pendant vertex of $T$ distinct from $x$ and $y$ must be dominated by $z$; otherwise, \{x,y,z\} would fail to be a dominating set, and the singleton set $\{z\}$ would not be able to form a secure coalition with any singleton part of $\pi$.

Let $x_1$ and $y_1$ denote the support vertices of $x$ and $y$, respectively. Since $\{x,y,z\}$ must be a dominating set, the vertex $z$ must be adjacent to both $x_1$ and $y_1$. Indeed, if $z\not\sim y_1$, then $\{x,y\}\cup\{y_1\}$ is not a dominating set. A similar argument applies to $x_1$.

Observe that $x_1$ and $y_1$ can dominate only the vertices $x$, $y$, and $z$. If either $x_1$ or $y_1$ dominates an additional vertex $w$, then $\{x,y\}\cup\{w\}$ is not a dominating set, and the singleton set $\{w\}$ cannot form a secure coalition with any singleton part of $\pi$. Therefore, $d(x_1)=d(y_1)=2$. It follows that every vertex in $V(T)\setminus\{x,y,x_1,y_1\}$
must be dominated by $z$. Since $T\not\cong P_5$, the vertex $z$ has a neighbor other than $x_1$ and $y_1$. One can now verify that neither $\{x_1\}$ nor $\{y_1\}$ can form a secure coalition with any set of $\pi$, contradicting the assumption that $\pi$ is a secure coalition partition.

Therefore, for every tree $T$ of order $n\ge 5$ with $T\not\cong P_5$, we have $SEC(T)\neq n-1$. Thus, it follows that
$SEC(T)\le n-2$.
\end{proof}

\end{document}
